\documentclass[12pt, twoside]{article}
%\documentclass[12pt, twoside]{article}
\usepackage[letterpaper, margin=1in, headsep=0.2in]{geometry}
\setlength{\headheight}{0.6in}
%\usepackage[english]{babel}
\usepackage[utf8]{inputenc}
\usepackage{microtype}
\usepackage{amsmath}
\usepackage{amssymb}
%\usepackage{amsfonts}
\usepackage[nomessages]{fp} %\FPeval{\var-name}{2*sin(pi/6)}
\usepackage{siunitx} %units in math. eg 20\milli\meter
\usepackage{yhmath} % for arcs, overparenth command
\usepackage{tikz} %graphics
\usetikzlibrary{quotes, angles, arrows, arrows.meta}
\usepackage{pgfplots}
\usepgfplotslibrary{fillbetween}
\pgfplotsset{compat=1.16}
\usepackage{graphicx} %consider setting \graphicspath{{images/}}
\usepackage{parskip} %no paragraph indent
\usepackage{enumitem}
\usepackage{multicol}
\usepackage{venndiagram}

\usepackage{fancyhdr}
\pagestyle{fancy}
\fancyhf{}
\renewcommand{\headrulewidth}{0pt} % disable the underline of the header
\raggedbottom
\hfuzz=2mm %suppresses overfull box warnings

\usepackage{hyperref}

\fancyhead[LO]{La Scuola d'Italia / Dr.\ Huson / IB Math 11 \\* 30 September 2026}
\fancyhead[RO]{\fbox{\begin{minipage}{6cm}\footnotesize Nickname:\\[0.5cm]\end{minipage}}}
\fancyhead[LE]{\fbox{\begin{minipage}{6cm}\footnotesize Nickname:\\[0.5cm]\end{minipage}}}
\fancyfoot[C]{\thepage}

\begin{document}
\subsubsection*{1.7 Classwork: Financial calculations and compound interest}

\begin{enumerate}[itemsep=0.15cm]

%--- Problem 1: Simple interest as an arithmetic sequence [6 marks] ---
%--- Source: Original (freshly authored, AA SL 1.2) ---
\item[1.]

With \textbf{simple interest}, the interest is paid only on the original
amount invested, so the same interest $I$ is added every year:
$$I = P \times \frac{r}{100},$$
where $P$ is the amount invested and $r$ is the annual interest rate
in percent.

Maria invests \$$2000$ in an account that pays simple interest at a
rate of $3.5\%$ per year.

\begin{enumerate}[itemsep=0.1cm]
  \item Find the interest Maria earns each year.
  \item Write down the value of her investment at the end of years
  $1$, $2$ and $3$.
  \item Write down an expression for the value $V_{n}$ of the
  investment at the end of $n$ years.
  \item Find the value of the investment at the end of $12$ years.
  \item Find the number of complete years after which the value of
  the investment first exceeds \$$3000$.
\end{enumerate}

\begin{tikzpicture}
  \draw (-0.6,0) rectangle (14.6,12.4);
  \foreach \y in {1,2,...,12} \draw[dotted] (0,\y)--(14.1,\y);
\end{tikzpicture}

\newpage

%--- Problem 2: Compound interest, annual vs monthly, doubling time [7 marks] ---
%--- Source: Original (freshly authored, AA SL 1.4) ---
\item[2.]

With \textbf{compound interest}, the interest is added to the account
and then earns interest itself. If \$$P$ is invested at $r\%$ per year,
compounded $k$ times per year, the value after $n$ years is
$$FV = P\left(1 + \frac{r}{100k}\right)^{kn}.$$

Luca invests \$$2000$ in Bank A, which pays $3.5\%$ per year
compounded \textbf{annually}.

\begin{enumerate}[itemsep=0.1cm]
  \item Find the value of Luca's investment at the end of $12$ years.
\end{enumerate}

Bank B pays $3.4\%$ per year compounded \textbf{monthly}.

\begin{enumerate}[itemsep=0.1cm]
  \item[(b)] Find the value of \$$2000$ invested in Bank B at the end
  of $12$ years.
  \item[(c)] State which bank gives the greater value after $12$ years.
  \item[(d)] For Bank A, find the number of complete years after which
  Luca's investment first becomes at least double its original value.
\end{enumerate}

\begin{tikzpicture}
  \draw (-0.6,0) rectangle (14.6,13.8);
  \foreach \y in {1,2,...,13} \draw[dotted] (0,\y)--(14.1,\y);
\end{tikzpicture}

\end{enumerate}
\end{document}
